% =====================================================================
% Capítulo 2 · Planificación
% Borrador generado a partir de docs/tfg/{PLAN,PROGRESO,BITACORA,DECISIONES}.md,
% AGENTS.md, .claude/, .github/workflows/ci.yml y el historial de git.
% Estado de las fuentes: 2026-09-29 (F2.4b cerrada).
% Paquetes necesarios: booktabs, tabularx, todonotes, hyperref.
% =====================================================================

\chapter{Planificación}
\label{cap:planificacion}

En este capítulo se describe cómo se ha organizado el trabajo: la metodología
seguida, la descomposición en fases y su relación con la Estructura de
Descomposición del Trabajo (EDT) de la propuesta, la temporización, la gestión
de riesgos, las herramientas empleadas y la estructura de costes.

% ---------------------------------------------------------------------
\section{Metodología}
\label{sec:metodologia}

\subsection{Desarrollo incremental por fases}

El proyecto se ha abordado con un modelo de desarrollo \textbf{incremental}:
el trabajo se divide en fases (F0--F8) y, dentro de las fases más grandes, en
incrementos más pequeños (por ejemplo, F2.1--F2.9 en la integración del CMS).
Cada fase o incremento tiene tres elementos fijados de antemano en el plan del
proyecto:

\begin{itemize}
  \item un \textbf{objetivo} y una lista de entregables;
  \item un \textbf{criterio de «hecho»} verificable (por ejemplo, «la
        aplicación arranca en local, CU-01 funciona a mano y la integración
        continua está en verde»);
  \item las \textbf{dependencias} con otras fases.
\end{itemize}

Un incremento solo se da por cerrado cuando supera su criterio de «hecho» y la
verificación común descrita en la sección~\ref{sec:verificacion-proceso}. Así,
al final de cada incremento el sistema queda en un estado funcional y
comprobado, y los retrasos afectan al alcance de los incrementos posteriores,
no a la estabilidad de lo ya construido. Esta forma de trabajar encaja con el
objetivo del TFG, una arquitectura \emph{reutilizable}: cada incremento añade
un módulo que debe poder integrarse sin romper los anteriores (RNF-04,
RNF-06).

Las decisiones de arquitectura o de alcance se registran en el momento en que
se toman como \textit{Architecture Decision Records} (ADR, registros de
decisiones de arquitectura)~\cite{nygard2011adr}, con su contexto, las
alternativas consideradas y sus consecuencias. En el momento de redactar este
borrador hay dieciséis (ADR-001 a ADR-016); los relevantes para el diseño se
resumen en el capítulo~\ref{cap:diseno}. Las incidencias (fallos propios o
heredados, problemas de entorno y de herramientas) se anotan también en el
momento en una bitácora con el formato \emph{qué pasó, causa, solución,
evidencia y lección}; esa bitácora es la fuente de los riesgos materializados
de la sección~\ref{sec:riesgos}.

\subsection{Desarrollo asistido por agentes de IA}
\label{sec:metodologia-ia}

% PENDIENTE (autor): comprobar si la ESI/UCA exige una declaración formal del
% uso de herramientas de IA generativa en el TFG (formato y ubicación) y
% ajustar este apartado a esa normativa.

Una parte sustancial del código, de las pruebas y de la documentación se ha
producido con la ayuda de un asistente de programación basado en modelos de
lenguaje que actúa como \textit{agente} (un programa que, a partir de una
instrucción, lee ficheros, ejecuta órdenes y propone o aplica cambios en el
repositorio). En concreto se ha usado Claude Code~\cite{claudecode}. Todos los
\textit{commits} del repositorio registran esa coautoría mediante la línea
\texttt{Co-Authored-By} del mensaje.
\todo{Confirmar el modelo o modelos de IA empleados y el periodo de uso para
la declaración de uso de IA.}

Para que este modo de trabajo sea controlable y evaluable, antes de escribir
código de la aplicación (fase F0) se construyó un \textit{harness}, es decir,
un arnés de directrices y controles automáticos que acota lo que hacen los
agentes y comprueba el resultado. Sus componentes son:

\begin{itemize}
  \item \textbf{Directrices escritas.} Un fichero raíz de directrices para
        agentes y desarrolladores (\texttt{AGENTS.md}) y otros por paquete que
        fijan el \textit{stack}, las convenciones, las reglas de seguridad de
        la base de datos, el idioma de comentarios y documentación (ADR-003) y
        una lista de verificación obligatoria al terminar cualquier cambio.
  \item \textbf{Documentación de seguimiento} en \texttt{docs/tfg/}: plan,
        requisitos, decisiones, trazabilidad, progreso, dedicación y bitácora.
        Un agente debe leer el progreso y el plan de la fase antes de empezar,
        y no puede reabrir una decisión sin registrar un ADR nuevo.
  \item \textbf{Procedimientos reutilizables} (\textit{skills}) para tareas
        repetitivas o delicadas: portar un módulo, reescribir comentarios,
        registrar una decisión, revisar un cambio de forma adversarial, auditar
        las políticas de seguridad de la base de datos o redactar un borrador
        de la memoria.
  \item \textbf{Controles automáticos} (\textit{hooks}) que se ejecutan sin
        intervención: al inicio de cada sesión se inyecta la fase actual y sus
        tareas pendientes; tras cada edición se comprueba que el fichero no
        contiene referencias prohibidas (ADR-007); y antes de terminar se
        recuerda actualizar el progreso si hay cambios de código sin
        registrar.
  \item \textbf{Permisos} explícitos: el código de partida es de solo lectura
        para los agentes y no pueden leer ficheros de secretos.
\end{itemize}

\paragraph{Reparto de responsabilidades.}
El reparto entre el autor y los agentes ha sido el siguiente:

\begin{itemize}
  \item \textbf{El autor} fija el alcance y los requisitos, toma las decisiones
        de arquitectura (por ejemplo, la integración del CMS en la consola de
        administración, ADR-011, fue una propuesta suya), revisa y acepta o
        rechaza los cambios, realiza pruebas manuales exploratorias y es el
        único que autoriza los \textit{commits} y su publicación. Las
        directrices prohíben a los agentes hacer \textit{commit} o
        \textit{push} sin su orden expresa.
  \item \textbf{Los agentes} analizan el código existente, proponen
        alternativas, implementan los incrementos, escriben las pruebas,
        ejecutan la verificación, redactan la documentación de seguimiento y
        los borradores de la memoria, y realizan las revisiones adversariales.
\end{itemize}
\todo{El autor debe revisar y completar este reparto con su propia valoración:
qué partes ha escrito o reescrito a mano, qué criterio ha seguido para aceptar
cambios y cuánto tiempo ha dedicado a la revisión.}

\paragraph{Cómo se verifica el trabajo.}
\label{sec:verificacion-proceso}
Ningún cambio se acepta por el mero hecho de haberlo producido un agente. Todo
cambio pasa por la misma cadena de comprobaciones, que el
\texttt{AGENTS.md} declara obligatoria:

\begin{enumerate}
  \item comprobación de tipos (TypeScript estricto), \textit{lint} y formato;
  \item comprobación de desmarcado (\texttt{check-branding});
  \item pruebas afectadas: unitarias, de base de datos (pgTAP) y
        \textit{end-to-end} (E2E) si se toca un flujo de usuario;
  \item revisión adversarial del \textit{diff} en dos pistas paralelas
        (calidad y arquitectura; búsqueda de errores y fallos de seguridad);
  \item si el cambio afecta a la base de datos, auditoría adversarial de las
        políticas de seguridad a nivel de fila con pruebas de ataque ejecutables
        y una etapa de refutación independiente (véase la
        sección~\ref{sec:verificacion-seguridad});
  \item revisión de los comentarios del código;
  \item actualización del progreso, la trazabilidad y, si procede, las
        decisiones y la bitácora.
\end{enumerate}

Las comprobaciones automatizables se repiten en la integración continua
(sección~\ref{sec:herramientas}). La revisión humana sigue siendo necesaria:
la incidencia B-14 (la consola de administración se cargaba sin segundo
factor) la detectó el autor con una prueba manual que ninguna prueba
automática cubría. De igual modo, la incidencia B-11 muestra que un test en
verde no garantiza la corrección: un test heredado afirmaba precisamente el
comportamiento inseguro. Por eso la metodología combina la verificación
automática con la lectura crítica de las aserciones y con pruebas
exploratorias del autor.

% ---------------------------------------------------------------------
\section{Fases y relación con la EDT}
\label{sec:fases}

La propuesta del TFG define una EDT con seis paquetes de trabajo:
(1)~Análisis, (2)~Arquitectura, (3)~Modelo de datos, (4)~Implementación,
(5)~Validación y (6)~Documentación. La tabla~\ref{tab:fases} relaciona cada
fase con esos paquetes.

\begin{table}[htbp]
  \centering
  \small
  \caption{Fases del proyecto, paquetes de la EDT que cubren y criterio de «hecho».}
  \label{tab:fases}
  \begin{tabularx}{\textwidth}{@{}l>{\raggedright\arraybackslash}p{3.4cm}cX@{}}
    \toprule
    \textbf{Fase} & \textbf{Nombre} & \textbf{EDT} & \textbf{Hecho cuando\ldots} \\
    \midrule
    F0 & Harness y planificación & 1 & El control de desmarcado pasa, los \textit{hooks} funcionan y existe el primer \textit{commit} \\
    F1 & Base SaaS importada y funcionando & 4 & La aplicación arranca en local, CU-01 funciona a mano y la CI está en verde \\
    F2 & Integración del CMS & 2, 3, 4 & Cada incremento F2.1--F2.9 cumple su criterio (BD, API, interfaz, explorador, usuarios y almacenamiento, auditoría, RBAC, paneles, cierre) \\
    F3 & Desmarcado completo e i18n en español & 4 & Control de desmarcado sin excepciones e interfaz en español por defecto \\
    F4 & Comentarios didácticos en español & 4, 6 & Sin comentarios en inglés en los módulos críticos y revisión aprobada \\
    F5 & Adaptación a pymes y configurabilidad & 2, 3, 4 & CU-02 a CU-05 funcionan y la configuración está documentada \\
    F6 & Validación & 5 & CI en verde y resultados documentados \\
    F7 & Despliegue & 4, 6 & Despliegue de demostración accesible y manual probado desde cero \\
    F8 & Memoria y manuales & 1, 2, 3, 6 & Continua; se cierra tras F6 y F7 \\
    \bottomrule
  \end{tabularx}
\end{table}

Las dependencias principales son lineales (F0 $\rightarrow$ F1 $\rightarrow$
F2 $\rightarrow$ F3 $\rightarrow$ F5 $\rightarrow$ F6), con dos excepciones:
F4 puede solaparse con F3 y F5, y F8 se desarrolla de forma continua, de modo
que cada fase cerrada deja un borrador de su parte de la memoria.

La fase F2 es la de mayor esfuerzo. Al comenzarla se evaluaron tres formas de
integrar el CMS (aplicaciones separadas, integración híbrida e integración
total) y se eligió la integración total en la consola de administración
(ADR-011), a pesar de su mayor coste estimado, por dar lugar a un único
servicio, un único inicio de sesión y un único estilo de código. Para
controlar ese riesgo, la fase se dividió en nueve incrementos y se identificó
de antemano el candidato a recortar si el plazo lo exige: los paneles
configurables (RF-11, requisito deseable).

\paragraph{Estado en la fecha de este borrador.}
F0 y F1 están cerradas. En F2 están cerrados los incrementos F2.1 (base de
datos), F2.2 (API), F2.3 (base de la interfaz), F2.4a (listado del explorador
de datos) y F2.4b (ficha de registro).
% PENDIENTE: actualizar este párrafo al integrar el capítulo.

% ---------------------------------------------------------------------
\section{Temporización}
\label{sec:temporizacion}

La dedicación se registra en el documento de progreso del proyecto. Las cifras
de la tabla~\ref{tab:dedicacion} son \textbf{estimaciones} obtenidas a partir
de la hora de los \textit{commits} (desde el inicio de la sesión hasta el
último \textit{commit} de cada tramo), redondeadas al cuarto de hora; el autor
debe revisarlas y registrar las siguientes sesiones con su dedicación real.

\begin{table}[htbp]
  \centering
  \small
  \caption{Dedicación registrada (estimada) por tramos de trabajo.}
  \label{tab:dedicacion}
  \begin{tabularx}{\textwidth}{@{}llXr@{}}
    \toprule
    \textbf{Fecha} & \textbf{Tramo} & \textbf{Fases} & \textbf{Horas} \\
    \midrule
    29/09/2026 & 11:15--11:58 & F0 (harness y plan) & $\approx$ 0,75 \\
    29/09/2026 & 11:58--12:35 & F1 (base SaaS, Docker, CI) & $\approx$ 0,75 \\
    29/09/2026 & 12:35--14:19 & F2 (análisis del CMS y decisión de integración) + F2.1 (BD y auditoría de seguridad) & $\approx$ 1,75 \\
    29/09/2026 & 14:19--15:05 & F2.2 (API del CMS) & $\approx$ 0,75 \\
    29/09/2026 & 15:05--15:50 & F2.3 (base de la interfaz) & $\approx$ 0,75 \\
    29/09/2026 & 15:50--17:02 & Corrección del MFA en la consola + F2.4a (listado) & $\approx$ 1,25 \\
    29/09/2026 & 17:02--17:50 & F2.4b (ficha de registro) + bitácora y trazabilidad & $\approx$ 0,75 \\
    \midrule
    \multicolumn{3}{@{}l}{\textbf{Total registrado}} & $\approx$ \textbf{6,75} \\
    \bottomrule
  \end{tabularx}
\end{table}

\todo{Completar la tabla con las sesiones posteriores y con la dedicación
real revisada por el autor (incluido el tiempo de revisión, pruebas manuales,
reuniones con el tutor y redacción que no deja \textit{commits}).}

\todo{Añadir la planificación inicial estimada por fase (horas previstas) para
compararla con la real, y un diagrama de Gantt (p.\,ej.\ con el paquete
\texttt{pgfgantt}) con las fechas previstas y reales de cada fase.}

% PENDIENTE (autor): las horas registradas corresponden a sesiones de trabajo
% asistido por agentes (sección~\ref{sec:metodologia-ia}). Conviene explicar en
% la memoria qué incluyen y qué no, para que la comparación con la carga de
% créditos del TFG sea honesta.

% ---------------------------------------------------------------------
\section{Gestión de riesgos}
\label{sec:riesgos}

La tabla~\ref{tab:riesgos} recoge los riesgos identificados. Para cada uno se
indica la medida de mitigación y si se ha llegado a materializar, con la
incidencia de la bitácora que lo documenta (identificadores B-xx).

\begin{table}[htbp]
  \centering
  \footnotesize
  \caption{Riesgos del proyecto, mitigación y materialización.}
  \label{tab:riesgos}
  \begin{tabularx}{\textwidth}{@{}l>{\raggedright\arraybackslash}p{3.6cm}>{\raggedright\arraybackslash}Xl@{}}
    \toprule
    \textbf{ID} & \textbf{Riesgo} & \textbf{Mitigación} & \textbf{Materializado} \\
    \midrule
    R-01 & Desbordamiento del plazo por el esfuerzo de integración del CMS &
      Incrementos F2.1--F2.9 funcionales; RF-11 identificado como primer recorte (ADR-011) & No (en seguimiento) \\
    R-02 & Fallos de seguridad en el código de partida &
      Auditoría adversarial de RLS con pruebas pgTAP y refutación independiente; pruebas de regresión (ADR-015) & Sí: B-05--B-09 \\
    R-03 & Exposición del entorno de desarrollo (servidor remoto con IP pública) &
      Filtrado en la cadena \texttt{DOCKER-USER}, acceso por túnel SSH y revisión de cada servicio nuevo (ADR-010) & Sí: B-04 \\
    R-04 & Errores introducidos por cambios masivos o por agentes de IA &
      Controles automáticos ampliados a las variantes del error, revisión adversarial y prueba manual del autor & Sí: B-01, B-14, B-16 \\
    R-05 & Deriva entre el esquema declarativo y las migraciones &
      \texttt{supabase db diff} en la verificación de todo cambio de BD & Sí: B-10 \\
    R-06 & Pruebas que no prueban lo que dicen &
      Leer las aserciones de las pruebas heredadas, no solo sus nombres & Sí: B-11 \\
    R-07 & Interrupción de las herramientas o de la conexión durante una tarea &
      Cerrar cada incremento con su \textit{commit}; tareas de seguridad formuladas como verificación defensiva y divididas en etapas & Sí: B-12, B-15 \\
    R-08 & Falsos positivos de las herramientas automáticas (secretos, pruebas inestables) &
      Documentar las credenciales públicas de desarrollo; ejecutar los E2E contra la \textit{build} de test & Sí: B-03, B-13 \\
    R-09 & Conflicto de licencia del código de partida &
      Repositorio privado y origen documentado internamente (ADR-001, ADR-006) & No \\
    R-10 & Validación de la reutilización sin un dominio de negocio concreto &
      Medir el arranque de un SaaS nuevo a partir de la plataforma (ADR-005, punto abierto P-01) & Pendiente \\
    \bottomrule
  \end{tabularx}
\end{table}

\todo{Asignar probabilidad e impacto (p.\,ej.\ bajo/medio/alto) a cada riesgo
según el criterio del autor y, si la plantilla lo pide, una matriz de
riesgos.}

Entre los riesgos materializados destacan tres por su enseñanza para la
planificación:

\begin{itemize}
  \item \textbf{B-04, exposición del entorno.} El entorno de desarrollo es un
        servidor remoto. Los servicios locales de la base de datos publicaban
        sus puertos en todas las interfaces y el motor de contenedores
        insertaba reglas que se saltaban el cortafuegos, de modo que la API y
        la base de datos (con credenciales de desarrollo conocidas) fueron
        accesibles desde Internet durante aproximadamente una hora, solo con
        datos de prueba. Se corrigió con una regla persistente, se recrearon
        los volúmenes y se pasó a acceder mediante un túnel SSH (ADR-010).
        Además, dio lugar a un criterio de revisión para cualquier servicio
        nuevo y a una advertencia que figurará en el manual de instalación.
  \item \textbf{B-05 a B-09, fallos heredados.} La auditoría de la base de
        datos del CMS encontró dos brechas confirmadas y varios
        endurecimientos necesarios en el código de partida (se analizan en la
        sección~\ref{sec:verificacion-seguridad}). El riesgo se había previsto
        y la mitigación (auditoría obligatoria en todo cambio de BD) fue la
        que permitió detectarlos.
  \item \textbf{B-15, corte de conexión con un agente trabajando.} El trabajo
        quedó a medias sin \textit{commit} y mezclado con una corrección de
        seguridad pendiente. La lección se incorporó al proceso: cada
        incremento se cierra con su propio \textit{commit}.
\end{itemize}

% ---------------------------------------------------------------------
\section{Herramientas}
\label{sec:herramientas}

\begin{itemize}
  \item \textbf{Control de versiones: git y GitHub.} Repositorio privado
        (condición de la licencia del código de partida, ADR-001). La rama
        principal es \texttt{main} y se trabaja en ramas por fase
        (\texttt{fase-1/base}, \texttt{fase-2/cms}). Los mensajes de
        \textit{commit} están en español y siguen la especificación
        \textit{Conventional Commits}~\cite{conventionalcommits}, con el
        módulo como ámbito.
  \item \textbf{Integración continua: GitHub Actions.} Un único
        \textit{workflow} con tres trabajos: (1)~controles propios del TFG
        (el \textit{scope} de paquetes es bloqueante; el desmarcado completo
        y el informe de comentarios en inglés son informativos hasta cerrar
        F3 y F4); (2)~comprobación de tipos, \textit{lint}, formato y pruebas
        unitarias; y (3)~base de datos local, pruebas pgTAP y E2E con
        Playwright contra la \textit{build} de test. Para ahorrar minutos de
        ejecución, el tercero solo se lanza en las \textit{pull requests}
        hacia \texttt{main} y a mano.
  \item \textbf{Monorepo:} Turborepo~\cite{turborepo} y pnpm, con un catálogo
        único de versiones.
  \item \textbf{Base de datos local:} Supabase CLI~\cite{supabase-docs} sobre
        Docker.
  \item \textbf{Pruebas:} Vitest (unitarias), pgTAP~\cite{pgtap} (base de
        datos) y Playwright~\cite{playwright} (E2E).
  \item \textbf{Asistente de programación:} Claude Code~\cite{claudecode},
        con el \textit{harness} descrito en la
        sección~\ref{sec:metodologia-ia}.
  \item \textbf{Scripts propios de control} (\texttt{scripts/tfg/}):
        comprobación de desmarcado, informe de comentarios en inglés y
        extracción de las etiquetas de trazabilidad del código.
  \item \textbf{Memoria:} \LaTeX{} sobre la plantilla de la ESI en Overleaf.
        % PENDIENTE (P-04): confirmar al integrar la plantilla oficial.
\end{itemize}

% ---------------------------------------------------------------------
\section{Costes}
\label{sec:costes}

% PENDIENTE (autor): todos los importes. No hay datos de costes en el
% repositorio; esta sección solo propone la estructura.

\todo{Completar la estimación de costes con datos reales.}

\begin{table}[htbp]
  \centering
  \small
  \caption{Estructura de costes del proyecto (pendiente de completar).}
  \label{tab:costes}
  \begin{tabularx}{\textwidth}{@{}Xrr@{}}
    \toprule
    \textbf{Concepto} & \textbf{Base de cálculo} & \textbf{Importe (€)} \\
    \midrule
    Personal: horas de desarrollo $\times$ coste/hora de un ingeniero junior & \textit{pendiente: horas y tarifa} & \textit{pendiente} \\
    Equipo de desarrollo (amortización) & \textit{pendiente} & \textit{pendiente} \\
    Servidor remoto de desarrollo (VPS) & \textit{pendiente: meses} & \textit{pendiente} \\
    Licencia comercial de la base de código de partida % PENDIENTE (P-06): cómo se presenta
      & \textit{pendiente} & \textit{pendiente} \\
    Suscripción al asistente de programación con IA & \textit{pendiente: meses} & \textit{pendiente} \\
    Servicios en la nube para el despliegue (F7) & \textit{pendiente} & \textit{pendiente} \\
    Otros (dominio, etc.) & \textit{pendiente} & \textit{pendiente} \\
    \midrule
    \textbf{Total} & & \textit{pendiente} \\
    \bottomrule
  \end{tabularx}
\end{table}

Las herramientas de desarrollo empleadas (git, Node.js, pnpm, Turborepo,
Supabase CLI, Docker, Vitest, pgTAP, Playwright) son de código abierto y no
tienen coste de licencia. Stripe se usa en modo de pruebas.
\todo{Confirmar el plan contratado de GitHub, del proyecto de Supabase en la
nube y de la plataforma de despliegue, y si generan coste.}
